\chapter{Introduction}
This chapter sets out the problem this phase addresses and the shape of the work that follows. Section~1.1 recalls the Phase~I system in one paragraph. Section~1.2 states the Phase~II problem. Section~1.3 lists the objectives of Phase~II, each with a short explanation. Section~1.4 defines the scope of the work across both phases. Section~1.5 maps the chapters that follow.

\section{Background}
A self-directed learner is a person who works through online courses, examination preparation, or self-study material without an instructor setting the pace. Such a learner typically begins with a fixed plan, of the form ``two hours a day, finish by this date.'' Within days the plan no longer matches reality, and the fixed finish date quietly turns wrong while the plan still displays it as though it were still valid. Phase~I of this project built an open-loop adaptive planner for this learner. The system estimated true study pace from session history with hierarchical Bayesian calibration, flagged genuine shifts in pace with CUSUM change-point detection, projected the completion date with a calibrated uncertainty band using Gaussian-process regression, and regenerated the schedule from realised capacity when the projection diverged from the deadline. Each component was compared against baselines on a synthetic dataset with known ground truth, using held-out evaluation with bootstrap confidence intervals and multiple-comparison correction, and the selected models were realised in a working local-first web application. The comparison reported its nulls honestly: the calibrated pace beat the fixed-rate baseline only after enough sessions had accumulated, and change detection stayed a documented robust null when no candidate dominated the frontier. What the open-loop system never checked was whether the logged hours produced learning.

\section{Problem Statement}
The pace signal in Phase~I comes from self-reported session logs. It records study time and says nothing about understanding. A learner can log every planned hour and still fail every test, while the plan keeps praising the pace and the examination result collapses. Existing tools leave the same gap open from the other side. Time-tracking applications record hours without judging understanding, and static plan generators never revise against any evidence at all. The missing piece is assessment-based verification of genuine learning, wired into the planning loop so that grades drive the decision to stay the course or replan. Phase~II closes this loop. It generates grounded assessments from the learner's own materials, grades them with an authoritative arm per question family, tracks mastery from the grade history, and feeds a verdict on the plan back into schedule regeneration.

\section{Objectives}
Phase~II of the project pursues the following objectives.
\begin{enumerate}
\item Generate grounded assessments in three families. Each generation call produces one question in a single family, objective, written, or coding, grounded in retrieved chunks of the learner's material with recorded citations. A mixed run alternates the family per call, so every question stays traceable to the material it tests.

\item Grade each family with its authoritative arm. Objective questions grade deterministically against authored answers, written answers grade against a rubric applied by a language model, and coding answers execute in a sandboxed runner with hidden tests. The server grade is authoritative. Any client-side check stays advisory and never reaches the record.

\item Estimate mastery with Bayesian knowledge tracing and steer difficulty from it. A Bayesian model updates a per-skill mastery probability from the graded response history, and the practice system adapts the difficulty of upcoming problems to the estimate instead of serving a fixed sequence.

\item Guide with a Socratic ladder and a gated reveal. The coach answers questions through a ladder of Nudge, Hint, Targeted explanation, and Reveal, and the final step stays locked until the learner confirms. Guidance therefore never hands over the answer before the learner has worked.

\item Return advisory feedback copy on each graded attempt. A server endpoint drafts learner-facing feedback grounded in the attempt and its citations for updated entries, with a static fallback when the model is unreachable. The copy never alters the grade.

\item Run mixed practice through the library and practice hub. Assessments and practice runs live under the material library with cache-first loading, and a practice run resolves its questions lazily from the loaded assessments. Starting practice therefore stays fast even on slow networks.

\item Replan from the loop verdict and validate the loop on synthetic longitudinal histories. A loop endpoint reads mastery, posture, and grade evidence and returns a hold or replan verdict under a fixed confidence gate, and a verdict for replan carries into schedule regeneration with its identity attached. Four synthetic learner archetypes over five weeks each supply the longitudinal histories that a single-semester project cannot collect from real learners.
\end{enumerate}

\section{Scope of the Project}
The work is delivered across two phases. \textbf{Phase~I}, reported earlier, built and evaluated the open-loop system: pace calibration, change detection, completion projection, and capacity-based schedule regeneration, each compared offline against baselines and realised in the application. \textbf{Phase~II}, reported here, closes the loop with assessment-based verification of genuine learning: grounded generation, authoritative grading, mastery tracing, Socratic guidance, advisory feedback, mixed practice runs, and verdict-driven replanning, validated on synthetic longitudinal histories. The target user in both phases is the individual self-directed learner. Multi-user, institutional, and instructor-mediated settings remain out of scope, as does deployment in a real classroom. The synthetic archetypes stand in for the longitudinal real-learner study a single semester cannot supply.

\section{Organization of the Report}
The remainder of the report is organized as follows. Chapter~2 surveys the research this phase draws on: Sections~2.1 to~2.4 cover knowledge tracing, item generation and tagging, grading, and guidance with feedback, Section~2.5 notes the sourcing method, and Section~2.6 summarizes the approaches and the gaps this project addresses. Chapter~3 specifies the system requirements in Sections~3.1 to~3.4. Chapter~4 describes the proposed methodology: the closed loop in Section~4.1, the selected models in Section~4.2, the comparison method in Section~4.3, the design evolution in Section~4.4, and the system architecture in Section~4.5. Chapter~5 gives the implementation details in Sections~5.1 to~5.5. Chapter~6 presents the results and discussion in Sections~6.1 to~6.5. Chapter~7 concludes in Section~7.1 and sets out future work in Section~7.2.
